\documentclass[a4paper,12pt]{extarticle}
\input{assessments/lecture04/quiz-style.tex}
\begin{document}
\quizheader{2 · NASM, BIOS и QEMU}
\input{assessments/lecture04/common-questions.tex}
\newpage
\section*{Вариант 2 · вопросы 7–10}
\begin{enumerate}
\setcounter{enumi}{6}
\setlength{\itemsep}{6mm}
\question{Какой размер имеет собранный загрузочный сектор из лабораторной?}
  {16 байтов.}{64 байта.}{512 байтов.}{1440 байтов.}
\question{Что записывается в одну позицию текстового VGA-экрана в учебном примере?}
  {Только числовая оценка.}{Код символа и байт атрибута.}
  {Адрес исходного файла и номер строки.}{Содержимое целого сектора диска.}
\question{Чем отличается вызов BIOS \texttt{INT 16h} от IRQ1 клавиатуры?}
  {Оба всегда инициирует инструкция нашей программы.}
  {\texttt{INT 16h} вызывает программа, а IRQ1 приходит от устройства.}
  {IRQ1 записывает оценки на диск, а \texttt{INT 16h} — в RAM.}
  {Различается только написание, действие одинаково.}
\question{Оценка появилась на экране, но сохранение на диск не выполнялось. Что будет после перезапуска учебной машины?}
  {Оценка сохранится благодаря видеопамяти.}
  {Оценка автоматически запишется в исходник NASM.}
  {Введённая оценка исчезнет: она оставалась в RAM.}
  {QEMU преобразует её в файл на диске.}
\end{enumerate}
\answergrid
\end{document}
